\subsection{\texorpdfstring{\(\mathbf{S}(\mathbf{M}^{*})\)、\(\mathbf{TS}(\mathbf{M})^{*\mathrm{gr}}\) 与 \(\operatorname{Pol}(\mathbf{M}, \mathbf{A})\) 之间的关系}{S(M*)、TS(M) 的分次对偶与 Pol(M, A) 之间的关系}}

设 M 为自由 A-模。我们将赋予分次对偶 \(\mathbf{TS}(\mathbf{M})^{*gr}\) 一个交换、结合且含单位元的分次 \(^{1}\) 代数结构，它由 \(\mathbf{TS}(\mathbf{M})\) 的分次余代数结构（III，第580页）导出。根据III，第497页，存在唯一的分次 A-代数同态

\[
\theta : \mathbf {S} (\mathbf {M} ^ {*}) \rightarrow \mathbf {T S} (\mathbf {M}) ^ {* \mathrm{gr}}
\]

它在次数 1 上诱导 \(\mathbf{M}^*\) 的恒等映射。

命题20. --- 若A-模M是自由且有限生成的，则 \(\theta\) 是分次代数的一个同构。

设 \((e_i)_{i\in I}\) 是 \(\mathbf{M}\) 的一组基，且 \((e_i^*)_{i\in I}\) 是 \(\mathbf{M}^*\) 的对偶基。对 \(\nu \in \mathbf{N}^I\) ，令

\[
e _ {\nu} = \prod_ {i \in I} \gamma_ {\nu_ {i}} (e _ {i}) \in \mathbf {T S} (\mathbf {M}).
\]

由命题4（IV，第47页），族 \((e_{\nu})_{\nu \in \mathbf{N}^{I}}\) 是 \(\mathbf{TS}(\mathbf{M})\) 的一组基；设 \((e_{\nu}^{*})\) 是 \(\mathbf{TS}(\mathbf{M})^{*\mathrm{gr}}\) 中与 \((e_{\nu})\) 对偶的基。根据III，第505页，定理1，只需证明对任意 \(\nu \in \mathbf{N}^{I}\) 有

\[
e _ {\nu} ^ {*} = \prod_ {i \in I} \left(e _ {i} ^ {*}\right) ^ {\nu_ {i}},
\]

或等价地，对 \(\rho, \sigma \in \mathbf{N}^{\mathrm{I}}\) 有 \(e_{\rho}^{*} \cdot e_{\sigma}^{*} = e_{\rho + \sigma}^{*}\) ；而最后一个断言可由IV，第50页，公式(12)得出。

注1. --- 同样地，我们可以看到，若 M 是有限生成的自由 A-模，则III，第593页所定义的分次代数 \(\mathbf{S}(\mathbf{M})^{*\mathrm{gr}}\) 可与 \(\mathbf{TS}(\mathbf{M}^*)\) 等同。

命题21. --- A-模的典范同态（IV，第55页）

\[
u: \mathbf {T S} (\mathrm{M}) ^ {* \mathrm{gr}} \rightarrow \operatorname{Pol} _ {\mathrm{A}} (\mathrm{M}, \mathrm{A})
\]

是一个代数同态。

设 \(a \in TS^{q}(M)^{*}\) , \(b \in TS^{r}(M)^{*}\) , \(x \in M\) ；我们有

\[
\begin{array}{r l} u (a b) (x) & = \langle a b, \gamma_ {q + r} (x) \rangle = \langle a \otimes b, c (\gamma_ {q + r} (x)) \rangle = \\ & = \langle a \otimes b, \gamma_ {q} (x) \otimes \gamma_ {r} (x) \rangle = \langle a, \gamma_ {q} (x) \rangle \langle b, \gamma_ {r} (x) \rangle = u (a) (x). u (b) (x), \end{array}
\]

由此即得结论。

注记。--- 2) 复合同态 \(\lambda_{\mathbf{M}}=u\circ\theta:\mathbf{S}(\mathbf{M}^{*})\to\mathrm{Pol}_{\mathbf{A}}(\mathbf{M},\mathbf{A})\) 是唯一的保单位代数同态，它诱导了

\[
\mathbf {M} ^ {*} = \operatorname{Pol} ^ {1} (\mathbf {M}, \mathbf {A})
\]

在 \(\operatorname{Pol}(\mathbf{M}, \mathbf{A})\) 中的包含映射。若 \(\mathbf{M}\) 是有限生成的自由模，且 \(\mathbf{A}\) 是无限整环，则 \(\lambda_{\mathbf{M}}\) 是双射的（命题20及命题19的推论）。特别地，若 \(\mathbf{A}\) 是无限整环，则典范同态 \(f \mapsto \tilde{f}\) 从 \(\mathbf{A}[\mathbf{X}_1, \ldots, \mathbf{X}_n]\) 到 \(\operatorname{Pol}(\mathbf{A}^n, \mathbf{A})\) （IV，第4页）是同构。

\begin{enumerate}
\def\labelenumi{\arabic{enumi})}
\setcounter{enumi}{2}
\tightlist
\item
  考虑余积 \(c_{\mathbf{S}}: \mathbf{S}(\mathbf{M}^{*}) \to \mathbf{S}(\mathbf{M}^{*} \times \mathbf{M}^{*})\) （III，第575页，例6）。对任意 \(v \in \mathbf{S}(\mathbf{M}^{*})\) , \(x, y \in \mathbf{M}\) ，多项式映射 \(\lambda_{\mathbf{M} \times \mathbf{M}}(c_{\mathbf{S}}(v)): \mathbf{M} \times \mathbf{M} \to \mathbf{A}\) 把 \((x, y)\) 映为 \(\lambda_{\mathbf{M}}(v)(x + y)\) 。如此定义的从 \(\mathbf{S}(\mathbf{M}^{*})\) 到 \(\operatorname{Map}(\mathbf{M} \times \mathbf{M}, \mathbf{A})\) 的两个代数同态在 \(\mathbf{M}^{*}\) 上一致，这依据关系式
\end{enumerate}

\[
(v \otimes 1 + 1 \otimes v) (x, y) = v (x + y) \quad (v \in \mathbf {M} ^ {*}).
\]

